\documentclass[12pt]{article}
\usepackage[utf8]{inputenc}
\usepackage[T1]{fontenc}
\usepackage{hyperref}
\title{Usku Language Resources for K-12 Curriculum Development}
\author{Compiled by Claude Sonnet 5 (claude-sonnet-5) for Andr\'as Kornai}
\date{}
\begin{document}
\maketitle
\section{Introduction}
Usku (also Afra; ISO 639-3 ulf) is a poorly documented Papuan language of Usku village, Senggi District, Keerom Regency, Papua, Indonesia. Sources give 20 or more speakers (Wikipedia) up to 110 (ASJP). Its classification is disputed: isolate, West Pauwasi, or unclassified. Very little exists.
\section{Primary Educational Resources}
None found.
\section{Linguistic and Reference Materials}
The best source is the SIL survey report by Im and Lebold (2006), \emph{Draft Survey Report on the Usku Language of Papua}, reportedly unpublished and not online: \url{https://glottolog.org/resource/reference/id/56786}. A short wordlist is in ASJP: \url{https://asjp.clld.org/languages/USKU_2}. Overview with basic vocabulary: \url{https://en.wikipedia.org/wiki/Usku_language}.
\section{Religious and Community Resources}
A people-group profile exists: \url{https://joshuaproject.net/people_groups/15746/ID}. No scripture or community organization was found.
\section{Digital and Audio Resources}
None found.
\section{Curriculum Development Context}
There are no school materials and no orthography was found. Any curriculum would first need documentation with the remaining speakers.
\end{document}